\documentclass[12pt]{article}
\usepackage{amsmath, amssymb, geometry, hyperref, enumitem}
\geometry{letterpaper, margin=1in}

\title{\textbf{Math - Test Sections Questions}\\
\large Section 1.2: Solutions and Initial Value Problems}
\author{}
\date{}

\begin{document}

\maketitle
\tableofcontents
\newpage

% ---------------------------------------------------------
\section{Concept Checklist \& Question Mapping}
% ---------------------------------------------------------
Before attempting the problem sets, ensure you understand the following concepts derived from the text. The questions are weighted and mapped as follows:

\begin{itemize}
    \item \textbf{Concept 1: Verification of Explicit Solutions \& Interval of Validity:} Testing if a given explicit function $y(x)$ satisfies the DE and determining where it is valid. (Questions 1, 2, 3, 4)
    \item \textbf{Concept 2: Verification of Implicit Solutions \& Algebraic Forcing:} Using implicit differentiation and algebraic manipulations to match an implicit equation to a DE. (Questions 5, 6, 7, 8)
    \item \textbf{Concept 3: Existence and Uniqueness (E\&U) Theorem:} Testing the continuity of $f(x,y)$ and $\frac{\partial f}{\partial y}$ at the initial condition to guarantee a unique solution. (Questions 9, 10, 11, 12)
    \item \textbf{Concept 4: Method of Undetermined Parameters:} Finding the constant parameter $m$ for generic families of test solutions like $y=x^m$ or $y=e^{mx}$. (Questions 13, 14, 15, 16)
    \item \textbf{Concept 5: Solving Initial Value Problems (IVPs):} Substituting initial conditions into general solutions to find the specific constant $C$. (Questions 17, 18, 19)
    \item \textbf{Concept 6: Singular Solutions \& Discontinuities:} Finding "lost" solutions not covered by the general family and evaluating IVPs centered on asymptotes/discontinuities. (Questions 20, 21, 22, 23)
    \item \textbf{Concept 7: Mathematical Modeling:} Translating real-world scenarios (finance, physics) into functional differential equations. (Questions 24, 25)
    \item \textbf{Concept 8: Advanced Diagnostic Testing ("Find the Flaw"):} Identifying and correcting specific calculus and conceptual errors in provided worked-out solutions. (Questions 26, 27, 28, 29, 30)
\end{itemize}

\newpage

% ---------------------------------------------------------
\section{Problem Sets}
% ---------------------------------------------------------

\subsection*{Part 1: Verification of Explicit Solutions}
\textbf{1.} Verify that $y = (8-2x)^{1/3}$ is an explicit solution to the differential equation $\frac{dy}{dx} = -\frac{2}{3y^2}$. State the intervals of validity for the general solution. \\
\textbf{2.} Show that $y = e^{3x} - x$ is a solution to the differential equation $\frac{dy}{dx} - 3y = 3x - 1$. \\
\textbf{3.} Determine whether $y = x^3 + \sin(x)$ is a solution to $x^2 y'' - 6y = -x^2 \sin(x) - 6\sin(x)$. \\
\textbf{4.} Determine whether $\theta = 2e^{-t} + e^{2t}$ is a solution to the differential equation $\frac{d^2\theta}{dt^2} - \theta = 3e^{2t}$.

\subsection*{Part 2: Verification of Implicit Solutions}
\textbf{5.} Verify that $x^4 + y^4 = 16$ is an implicit solution to the differential equation $\frac{dy}{dx} = -\frac{x^3}{y^3}$. Determine the interval restrictions for $x$. \\
\textbf{6.} Show that $y \cos(x) - x \sin(y) = C$ is an implicit solution to $\frac{dy}{dx} = \frac{y \sin(x) + \sin(y)}{\cos(x) - x \cos(y)}$. \\
\textbf{7.} Verify that $e^{xy} = y$ is an implicit solution to the differential equation $\frac{dy}{dx} = \frac{y^2}{1 - xy^2}$. \\
\textbf{8.} Verify that $x^2 y - y^2 x = 4$ is an implicit solution to $\frac{dy}{dx} = \frac{y^2 - 2xy}{x^2 - 2xy}$.

\subsection*{Part 3: Existence and Uniqueness Theorem}
\textit{For the following problems, determine if the Existence and Uniqueness Theorem guarantees a unique solution for the given Initial Value Problem (IVP).} \\
\textbf{9.} $\frac{dy}{dx} = 4x^2 y^3$ with $y(0) = 2$. \\
\textbf{10.} $\frac{dy}{dx} = \sqrt{y - 1}$ with $y(1) = 1$. \\
\textbf{11.} $\frac{dy}{dx} = \frac{x}{y^2 - 4}$ with $y(0) = 2$. \\
\textbf{12.} $\frac{dy}{dx} = \frac{x^2}{y + 1}$ with $y(2) = 0$.

\subsection*{Part 4: Method of Undetermined Parameters}
\textbf{13.} Find the value(s) of the parameter $m$ such that $y = x^m$ is a solution to $x^2 y'' - 5xy' + 8y = 0$. \\
\textbf{14.} Find the value(s) of the parameter $m$ such that $y = x^m$ is a solution to $x^2 y'' + 3xy' - 3y = 0$. \\
\textbf{15.} Find the value(s) of the parameter $m$ such that $y = e^{mx}$ is a solution to $y'' - 4y' - 5y = 0$. \\
\textbf{16.} Find the value(s) of the parameter $m$ such that $y = e^{mx}$ is a solution to $y'' + 6y' + 9y = 0$.

\subsection*{Part 5: Solving Initial Value Problems}
\textbf{17.} Given the general solution $y = C e^{-2x} + \frac{1}{2}x$, find the particular solution satisfying $y(0) = 3$. \\
\textbf{18.} Given that $x^2 + 2y^2 = C$ is a general implicit solution to a DE, find the particular solution that satisfies $y(2) = 1$. \\
\textbf{19.} Given the general solution $y = \frac{C}{x-1}$, find the particular solution satisfying $y(2) = 4$.

\subsection*{Part 6: Singular Solutions \& Discontinuities}
\textbf{20.} The DE $\frac{dy}{dx} = 2\sqrt{y}$ has a general one-parameter family of solutions $y = (x+C)^2$. Show that $y = 0$ is a singular solution that cannot be obtained from the general family. \\
\textbf{21.} The DE $\frac{dy}{dx} = 3y^{2/3}$ has a general family of solutions $y = (x-C)^3$. Identify a singular solution and prove it cannot be obtained from the family. \\
\textbf{22.} Consider the DE $x \frac{dy}{dx} - 3y = 0$ with the general solution $y = C x^3$. Show that the IVP with $y(0) = 5$ has no solution, and explain why this does not contradict the E\&U theorem. \\
\textbf{23.} Using the same DE and general solution from Question 22, show that the IVP with $y(0) = 0$ has infinitely many solutions.

\subsection*{Part 7: Mathematical Modeling}
\textbf{24.} An investor's retirement account starts with \$500,000. It generates a continuous fixed return of 6\% per year. Simultaneously, the investor withdraws money continuously at a rate of \$40,000 per year for living expenses. Set up the IVP modeling the total value of the account $V(t)$ over time. \\
\textbf{25.} According to Newton's Law of Cooling, a cup of coffee initially at $90^\circ\text{C}$ is placed in a room kept at a constant $20^\circ\text{C}$. The rate of cooling is proportional to the difference between the coffee's temperature and the room's temperature, with a proportionality constant $k = 0.05 \text{ min}^{-1}$. Set up the IVP for the temperature $T(t)$.

\subsection*{Part 8: Advanced Diagnostic Testing ("Find the Flaw")}
\textit{Identify the specific mathematical error in the following worked-out solutions, explain why it is wrong, and provide the correct steps.} \\

\noindent \textbf{26. Flawed Problem 1:} Verify that $x^2 + 4y^2 = 36$ is a solution to $y' = -\frac{x}{4y}$. \\
\textbf{Flawed Solution:} Differentiating both sides with respect to $x$:
\begin{align*}
2x + 8y \cdot y' &= 36 \\
8y \cdot y' &= 36 - 2x \\
y' &= \frac{36 - 2x}{8y} = \frac{18 - x}{4y}
\end{align*}
Since this does not match $-\frac{x}{4y}$, it is not a solution. \\

\noindent \textbf{27. Flawed Problem 2:} Use the E\&U theorem to determine if $\frac{dy}{dx} = y^{3/4}$ with $y(0) = 0$ has a unique solution. \\
\textbf{Flawed Solution:} Here, $f(x,y) = y^{3/4}$. The initial point is $(0,0)$. Because the function $f(x,y) = y^{3/4}$ is continuous at $y = 0$, it is continuous in a rectangle around $(0,0)$. Therefore, the theorem guarantees a unique solution. \\

\noindent \textbf{28. Flawed Problem 3:} Determine if $y = e^{-x^2}$ is a solution to $y'' + 2y = 0$. \\
\textbf{Flawed Solution:} Find the first and second derivatives:
\begin{align*}
y &= e^{-x^2} \\
y' &= -2x e^{-x^2} \\
y'' &= (-2x)(-2x e^{-x^2}) = 4x^2 e^{-x^2} \quad \text{(applying chain rule again)}
\end{align*}
Substitute into DE: $4x^2 e^{-x^2} + 2(e^{-x^2}) \neq 0$. Not a solution. \\

\noindent \textbf{29. Flawed Problem 4:} Find $m$ such that $y = e^{mx}$ is a solution to $y'' - 7y' = 0$. \\
\textbf{Flawed Solution:} Let $y = e^{mx}$. Then $y' = m e^{mx}$ and $y'' = m^2 e^{mx}$. Substitute:
\begin{align*}
m^2 e^{mx} - 7m e^{mx} &= 0 \\
e^{mx}(m^2 - 7m) &= 0
\end{align*}
Dividing out $m e^{mx}$ leaves $m - 7 = 0$, so $m = 7$. \\

\noindent \textbf{30. Flawed Problem 5:} Determine the interval of validity for the solution $y = \frac{1}{x-2}$ to a DE, with initial condition $y(4) = \frac{1}{2}$. \\
\textbf{Flawed Solution:} The function $y = \frac{1}{x-2}$ and its derivative $y' = \frac{-1}{(x-2)^2}$ are both continuous and defined for all real numbers except $x = 2$. Therefore, the interval of validity is $(-\infty, 2) \cup (2, \infty)$.

\newpage

% ---------------------------------------------------------
\section{Detailed Solutions}
% ---------------------------------------------------------

\subsection*{Part 1: Verification of Explicit Solutions}
\textbf{Solution 1:} First, convert the fractional exponent and differentiate using the chain rule:
$$y = (8-2x)^{1/3}$$
$$\frac{dy}{dx} = \frac{1}{3}(8-2x)^{-2/3} \cdot (-2) = \frac{-2}{3(8-2x)^{2/3}}$$
Since $y = (8-2x)^{1/3}$, squaring it gives $y^2 = (8-2x)^{2/3}$. Substituting this back in:
$$\frac{dy}{dx} = -\frac{2}{3y^2}$$
This matches the DE. To find the interval of validity, the denominator $y^2$ cannot equal $0$. Thus, $8-2x \neq 0 \implies x \neq 4$. The intervals of validity are $(-\infty, 4)$ or $(4, \infty)$. \\

\noindent \textbf{Solution 2:} Find the first derivative of $y$:
$$y' = 3e^{3x} - 1$$
Substitute $y$ and $y'$ into the left side of the DE:
$$\text{LHS} = y' - 3y = (3e^{3x} - 1) - 3(e^{3x} - x)$$
$$= 3e^{3x} - 1 - 3e^{3x} + 3x = 3x - 1 = \text{RHS}$$
It is a valid solution. \\

\noindent \textbf{Solution 3:} Find the first and second derivatives of $y$:
$$y' = 3x^2 + \cos(x)$$
$$y'' = 6x - \sin(x)$$
Substitute into the left side of the DE:
$$\text{LHS} = x^2(6x - \sin(x)) - 6(x^3 + \sin(x))$$
$$= 6x^3 - x^2\sin(x) - 6x^3 - 6\sin(x) = -x^2\sin(x) - 6\sin(x) = \text{RHS}$$
It is a valid solution. \\

\noindent \textbf{Solution 4:} Find the first and second derivatives of $\theta$:
$$\frac{d\theta}{dt} = -2e^{-t} + 2e^{2t}$$
$$\frac{d^2\theta}{dt^2} = 2e^{-t} + 4e^{2t}$$
Substitute into the left side of the DE:
$$\text{LHS} = \frac{d^2\theta}{dt^2} - \theta = (2e^{-t} + 4e^{2t}) - (2e^{-t} + e^{2t}) = 3e^{2t} = \text{RHS}$$
It is a valid solution.

\subsection*{Part 2: Verification of Implicit Solutions}
\textbf{Solution 5:} Differentiate both sides with respect to $x$:
$$\frac{d}{dx}[x^4 + y^4] = \frac{d}{dx}[16]$$
$$4x^3 + 4y^3 \frac{dy}{dx} = 0$$
$$\frac{dy}{dx} = -\frac{4x^3}{4y^3} = -\frac{x^3}{y^3}$$
This matches the DE. To find the interval restriction, notice that $\frac{dy}{dx}$ is undefined when $y=0$. From the original equation, if $y=0$, then $x^4 = 16 \implies x = \pm 2$. Since $y$ must also be a real-valued function, $x^4 < 16$ for $y \neq 0$. Therefore, the interval is $x \in (-2, 2)$. \\

\noindent \textbf{Solution 6:} Differentiate both sides implicitly with respect to $x$ using the product rule:
$$\frac{d}{dx}[y \cos(x) - x \sin(y)] = \frac{d}{dx}[C]$$
$$\left( y' \cos(x) - y \sin(x) \right) - \left( \sin(y) + x \cos(y) y' \right) = 0$$
Group the $y'$ terms together:
$$y' \cos(x) - x \cos(y) y' = y \sin(x) + \sin(y)$$
$$y' (\cos(x) - x \cos(y)) = y \sin(x) + \sin(y)$$
$$y' = \frac{y \sin(x) + \sin(y)}{\cos(x) - x \cos(y)}$$
Matches the DE perfectly. \\

\noindent \textbf{Solution 7:} Differentiate both sides implicitly using the product and chain rules:
$$e^{xy} \left( 1 \cdot y + x y' \right) = y'$$
Substitute the original relationship $e^{xy} = y$ back into the derivative to simplify:
$$y(y + x y') = y'$$
$$y^2 + xy^2 y' = y'$$
$$y^2 = y' - xy^2 y'$$
$$y^2 = y'(1 - xy^2) \implies y' = \frac{y^2}{1 - xy^2}$$
Matches the DE perfectly. \\

\noindent \textbf{Solution 8:} Differentiate both sides implicitly with respect to $x$:
$$\frac{d}{dx}[x^2 y - y^2 x] = 0$$
$$(2x y + x^2 y') - (2y y' x + y^2) = 0$$
$$y'(x^2 - 2xy) + 2xy - y^2 = 0$$
$$y'(x^2 - 2xy) = y^2 - 2xy \implies y' = \frac{y^2 - 2xy}{x^2 - 2xy}$$
Matches the DE perfectly.

\subsection*{Part 3: Existence and Uniqueness Theorem}
\textbf{Solution 9:} Identify $f(x,y) = 4x^2 y^3$. The partial derivative with respect to $y$ is $\frac{\partial f}{\partial y} = 12x^2 y^2$. Both functions are polynomials, meaning they are continuous everywhere on the $xy$-plane. Therefore, a unique solution is \textbf{guaranteed} around $(0,2)$. \\

\noindent \textbf{Solution 10:} Identify $f(x,y) = \sqrt{y - 1}$. The partial derivative is $\frac{\partial f}{\partial y} = \frac{1}{2\sqrt{y - 1}}$. At the initial point $(1,1)$, the denominator becomes zero, meaning the partial derivative is discontinuous at $y=1$. Therefore, a unique solution is \textbf{not guaranteed}. \\

\noindent \textbf{Solution 11:} Identify $f(x,y) = \frac{x}{y^2 - 4}$. The initial condition is $(0,2)$. At $y=2$, the denominator of $f(x,y)$ equals zero, meaning the function itself is discontinuous at the initial point. Therefore, a unique solution is \textbf{not guaranteed}. \\

\noindent \textbf{Solution 12:} Identify $f(x,y) = \frac{x^2}{y + 1}$ and compute the partial derivative: $\frac{\partial f}{\partial y} = -\frac{x^2}{(y+1)^2}$. Evaluating both at the initial point $(2,0)$ yields $f(2,0) = 4$ and $\frac{\partial f}{\partial y}(2,0) = -4$. Since neither has division by zero or any other discontinuity at $(2,0)$, both are continuous in a region around it. A unique solution is \textbf{guaranteed}.

\subsection*{Part 4: Method of Undetermined Parameters}
\textbf{Solution 13:} Let $y = x^m$. Then $y' = m x^{m-1}$ and $y'' = m(m-1)x^{m-2}$. Substitute into the DE:
$$x^2(m(m-1)x^{m-2}) - 5x(m x^{m-1}) + 8(x^m) = 0$$
$$x^m [m^2 - m - 5m + 8] = 0$$
Assuming $x \neq 0$, solve the characteristic polynomial:
$$m^2 - 6m + 8 = 0 \implies (m-2)(m-4) = 0 \implies \mathbf{m = 2, 4}$$

\noindent \textbf{Solution 14:} Following the same setup for $y = x^m$:
$$x^m [m(m-1) + 3m - 3] = 0$$
$$m^2 - m + 3m - 3 = 0 \implies m^2 + 2m - 3 = 0$$
$$(m+3)(m-1) = 0 \implies \mathbf{m = -3, 1}$$

\noindent \textbf{Solution 15:} Let $y = e^{mx}$. Then $y' = m e^{mx}$ and $y'' = m^2 e^{mx}$. Substitute into the DE:
$$m^2 e^{mx} - 4m e^{mx} - 5 e^{mx} = 0$$
$$e^{mx}(m^2 - 4m - 5) = 0$$
Since $e^{mx} \neq 0$, solve:
$$m^2 - 4m - 5 = 0 \implies (m-5)(m+1) = 0 \implies \mathbf{m = 5, -1}$$

\noindent \textbf{Solution 16:} Following the same setup for $y = e^{mx}$:
$$e^{mx}(m^2 + 6m + 9) = 0$$
$$m^2 + 6m + 9 = 0 \implies (m+3)^2 = 0 \implies \mathbf{m = -3}$$

\subsection*{Part 5: Solving Initial Value Problems}
\textbf{Solution 17:} Substitute $x = 0$ and $y = 3$ into the general solution:
$$3 = C e^{-2(0)} + \frac{1}{2}(0)$$
$$3 = C(1) + 0 \implies C = 3$$
The particular solution is $y = 3e^{-2x} + \frac{1}{2}x$. \\

\noindent \textbf{Solution 18:} Substitute $x = 2$ and $y = 1$ into the implicit general solution:
$$(2)^2 + 2(1)^2 = C \implies 4 + 2 = C \implies C = 6$$
The particular solution is $x^2 + 2y^2 = 6$. \\

\noindent \textbf{Solution 19:} Substitute $x = 2$ and $y = 4$ into the general solution:
$$4 = \frac{C}{2 - 1} \implies 4 = \frac{C}{1} \implies C = 4$$
The particular solution is $y = \frac{4}{x-1}$.

\subsection*{Part 6: Singular Solutions \& Discontinuities}
\textbf{Solution 20:} First, check if $y=0$ satisfies the DE: $\frac{d}{dx}(0) = 2\sqrt{0} \implies 0 = 0$. It is a valid solution. Next, see if it can be derived from the family $y = (x+C)^2$. Setting $(x+C)^2 = 0$ implies $x = -C$. Because $C$ is a constant, it cannot equal $-x$ for all values of $x$. Thus, it is impossible for the general solution to yield the line $y=0$ everywhere. $y=0$ is a singular solution. \\

\noindent \textbf{Solution 21:} Check if $y=0$ satisfies the DE: $\frac{d}{dx}(0) = 3(0)^{2/3} \implies 0=0$. It is a solution. Set the general family equal to 0: $(x-C)^3 = 0 \implies x = C$. Just as in the previous problem, $C$ cannot be a variable. $y=0$ is a singular solution. \\

\noindent \textbf{Solution 22:} Substitute the initial condition $(0,5)$ into the general solution $y = C x^3$:
$$5 = C(0)^3 \implies 5 = 0$$
This is a mathematical contradiction, meaning no solution exists. To see why the E\&U theorem does not contradict this, isolate $\frac{dy}{dx}$ to define $f(x,y)$:
$$\frac{dy}{dx} = \frac{3y}{x}$$
Because $f(x,y)$ involves division by $x$, it is discontinuous at $x=0$. The theorem’s prerequisites fail, so uniqueness (or existence) is not guaranteed. \\

\noindent \textbf{Solution 23:} Substitute the initial condition $(0,0)$ into the general solution $y = C x^3$:
$$0 = C(0)^3 \implies 0 = 0$$
This identity holds true for \textit{any} real number $C$. Thus, there are infinitely many solutions (e.g., $y=x^3$, $y=5x^3$, $y=-2x^3$). Again, because $f(x,y) = \frac{3y}{x}$ is discontinuous at $x=0$, the E\&U theorem is silent, permitting multiple solutions to exist at that point.

\subsection*{Part 7: Mathematical Modeling}
\textbf{Solution 24:} The rate of change of the portfolio $\frac{dV}{dt}$ is equal to the money gained from interest minus the money removed via withdrawals.
Interest gained $= 0.06 V(t)$. Money removed $= 40,000$.
$$\frac{dV}{dt} = 0.06V - 40000$$
The initial condition is the starting balance: $V(0) = 500,000$. \\

\noindent \textbf{Solution 25:} Let $T(t)$ be the temperature of the coffee. The ambient temperature is $T_a = 20$. The rate of change is proportional (and negative, since it's cooling) to $(T - T_a)$.
$$\frac{dT}{dt} = -0.05(T - 20)$$
The initial condition is the starting temperature: $T(0) = 90$.

\subsection*{Part 8: Advanced Diagnostic Testing ("Find the Flaw")}
\textbf{Solution 26 (Flaw in Derivative of Constant):} The student correctly applied implicit differentiation to the variables but forgot to differentiate the constant $36$ on the right side. The derivative of a constant is 0, not 36.
\textit{Correct Steps:}
$$2x + 8y y' = 0 \implies 8y y' = -2x \implies y' = -\frac{2x}{8y} = -\frac{x}{4y}$$

\noindent \textbf{Solution 27 (Flaw in E\&U Partial Derivative):} The student only checked the continuity of $f(x,y)$. They failed to check the partial derivative $\frac{\partial f}{\partial y} = \frac{3}{4}y^{-1/4} = \frac{3}{4 y^{1/4}}$. At $y=0$, the partial derivative requires division by zero and is therefore discontinuous. The theorem does \textit{not} guarantee a unique solution. \\

\noindent \textbf{Solution 28 (Flaw in the Second Derivative):} The student improperly applied the chain rule to a product. The first derivative $y' = -2x e^{-x^2}$ is a product of two functions of $x$: $(-2x)$ and $(e^{-x^2})$. Calculating the second derivative requires the Product Rule.
\textit{Correct Steps:}
$$y'' = \frac{d}{dx}[-2x](e^{-x^2}) + (-2x)\frac{d}{dx}[e^{-x^2}] = -2e^{-x^2} + (-2x)(-2x e^{-x^2}) = -2e^{-x^2} + 4x^2 e^{-x^2}$$

\noindent \textbf{Solution 29 (Flaw in Factoring/Roots):} By dividing the entire equation by the variable $m$, the student improperly discarded a valid mathematical root. You can divide by $e^{mx}$ because it is strictly positive, but you must factor $m$ out.
\textit{Correct Steps:}
$$e^{mx} m (m - 7) = 0 \implies m = 0, \text{ or } m = 7$$

\noindent \textbf{Solution 30 (Flaw in Interval of Validity Definition):} The student wrote the domain of the function, not the interval of validity for the IVP. An interval of validity must be a single, unbroken continuous interval that contains the initial point. Since the initial point is $x=4$, we must pick the piece of the domain that contains 4 without crossing the asymptote at $x=2$.
\textit{Correct Answer:} The interval of validity is $(2, \infty)$.

\newpage

% ---------------------------------------------------------
\section{Concept Checklist \& Question Mapping (Appendix)}
% ---------------------------------------------------------
\textit{Use this appendix to easily cross-reference problems to review specific concepts before an exam.}

\begin{itemize}
    \item \textbf{Concept 1: Verification of Explicit Solutions \& Interval of Validity} $\implies$ Questions 1, 2, 3, 4
    \item \textbf{Concept 2: Verification of Implicit Solutions \& Algebraic Forcing} $\implies$ Questions 5, 6, 7, 8
    \item \textbf{Concept 3: Existence and Uniqueness (E\&U) Theorem} $\implies$ Questions 9, 10, 11, 12
    \item \textbf{Concept 4: Method of Undetermined Parameters} $\implies$ Questions 13, 14, 15, 16
    \item \textbf{Concept 5: Solving Initial Value Problems (IVPs)} $\implies$ Questions 17, 18, 19
    \item \textbf{Concept 6: Singular Solutions \& Discontinuities} $\implies$ Questions 20, 21, 22, 23
    \item \textbf{Concept 7: Mathematical Modeling} $\implies$ Questions 24, 25
    \item \textbf{Concept 8: Advanced Diagnostic Testing ("Find the Flaw")} $\implies$ Questions 26, 27, 28, 29, 30
\end{itemize}

\end{document}